\documentclass[11pt,a4paper]{article}
\usepackage{fontspec,amsmath,amssymb,amsthm,microtype}
\setmainfont{lmroman10-regular.otf}[BoldFont=lmroman10-bold.otf,ItalicFont=lmroman10-italic.otf,BoldItalicFont=lmroman10-bolditalic.otf]
\usepackage[margin=25mm]{geometry}
\usepackage[unicode,hidelinks]{hyperref}
\hypersetup{pdfauthor={LangevinLean Project}}
\setlength{\emergencystretch}{3em}
\setlength{\parskip}{0.4em}
\setlength{\parindent}{0pt}
\allowdisplaybreaks[2]
\raggedbottom
\title{\texorpdfstring{Conditional parameter-entropy estimates}{Conditional parameter-entropy estimates}}
\author{LangevinLean Project}
\date{}
\begin{document}\maketitle
Draft mathematical exposition; these notes may contain mistakes.


The conditional joint entropy theorem and data processing give finite parameter divergence and its numerical estimate. The lifted-target logarithmic Sobolev assumption and all interpolation regularity assumptions are retained. The initial term is the full initial joint entropy, and an estimate normalized by particle count is multiplied by that count before it bounds the unnormalized parameter divergence.


\section{\texorpdfstring{Finiteness of the parameter divergence}{Finiteness of the parameter divergence}}


The Gaussian transition makes \(\rho_k\) a probability, and the logarithmic Sobolev assumption normalizes \(\Pi_N\). Under the interpolation assumptions, the finite joint entropy conclusion holds for \(\gamma\ge0\). Data processing therefore gives \(\operatorname{KL}(p_{\theta\#}\rho_k\Vert\Pi_N^\theta)<\infty\).



This conclusion is first an extended-valued finiteness statement. It establishes that the projected divergence is an ordinary finite number before a numerical upper bound is applied. No extra parameter-density formula is required.




Earlier results used here: \cite[Section~3.1.3]{LangevinLean2026ULALSI}; \cite[Section~3.34.1]{LangevinLean2026ULALSI}; \cite[Section~3.1.1]{LangevinLean2026IPLA}; \cite[Section~3.1.6]{LangevinLean2026IPLA}; \cite[Section~3.3.2]{LangevinLean2026IPLA}.

\section{\texorpdfstring{The parameter estimate from the full joint bound}{The parameter estimate from the full joint bound}}


Assume \(N>0\), \(\beta>0\), \(\beta\le NL\), and \(0<\gamma\le\beta/(4NL^2)\), together with the conditional interpolation hypotheses. The joint theorem gives \(H_k\le e^{-\beta\gamma k/N}H_0+8\gamma ND L^2/\beta\). The joint divergence is finite, so data processing yields

\[\operatorname{KL}(p_{\theta\#}\rho_k\Vert\Pi_N^\theta)\le e^{-\beta\gamma k/N}H_0+\frac{8\gamma ND L^2}{\beta}.\]

Here \(H_0=\operatorname{KL}(\mu_0\Vert\Pi_N)\). At the same iteration \(k\), projection decreases only the divergence on the left; substitution of the joint bound leaves its initial quantity, exponential, and error term unchanged. The initial quantity includes parameter and latent information. The estimate compares parameter probability laws and does not itself give concentration about a likelihood maximizer.




Earlier results used here: \cite[Section~3.1.3]{LangevinLean2026ULALSI}; \cite[Section~3.34.1]{LangevinLean2026ULALSI}; \cite[Section~3.1.1]{LangevinLean2026IPLA}; \cite[Section~3.1.6]{LangevinLean2026IPLA}; \cite[Section~3.3.2]{LangevinLean2026IPLA}.

\section{\texorpdfstring{Parameter divergence under a linearly scaled lifted constant}{Parameter divergence under a linearly scaled lifted constant}}


Suppose the logarithmic Sobolev constant of the lifted target is \(N\alpha\). The conditional normalized theorem gives

\[\frac{H_k}{N}\le e^{-\alpha\gamma k}\frac{H_0}{N}+\frac{8\gamma L^2}{\alpha}\left(d_x+\frac{d_\theta}{N}\right).\]

To bound the unnormalized parameter divergence, multiply the complete upper bound by positive \(N\). Data processing gives upper bound

\[N\left[e^{-\alpha\gamma k}\frac{H_0}{N}+\frac{8\gamma L^2}{\alpha}\left(d_x+\frac{d_\theta}{N}\right)\right].\]

For constant \(N\alpha\) the multiplication by \(N\) is explicitly

\[\operatorname{KL}(p_\#\rho_k\Vert p_\#\Pi_N)\le N\left[e^{-\alpha\gamma k}\frac{H_0}{N}+\frac{8\gamma L^2}{\alpha}(d_x+d_\theta/N)\right].\]

Probability normalization and finite joint entropy justify the finite data-processing inequality. Multiplication by \(N\) restores the initial quantity \(H_0\) and gives error term \(8\gamma L^2D/\alpha\), exactly as direct substitution of \(\beta=N\alpha\) into the full estimate does. The interpolation and lifted logarithmic Sobolev hypotheses remain in force, and this conclusion concerns divergence between probability laws.




Earlier results used here: \cite[Section~3.1.3]{LangevinLean2026ULALSI}; \cite[Section~3.34.1]{LangevinLean2026ULALSI}; \cite[Section~3.1.1]{LangevinLean2026IPLA}; \cite[Section~3.1.6]{LangevinLean2026IPLA}; \cite[Section~3.3.2]{LangevinLean2026IPLA}.

\begin{thebibliography}{ULA-LSI}
\bibitem[ULA-LSI]{LangevinLean2026ULALSI}
LangevinLean Project. \emph{Unadjusted Langevin under LSI}. LL-02, 2026. Draft report, content revision d01d1d7b45e1. \url{https://odakyildiz.github.io/LangevinLean/docs/pdf/UnadjustedLangevin/LSI/FullProof.pdf}.
\bibitem[IPLA-SC]{LangevinLean2026IPLA}
LangevinLean Project. \emph{Interacting Particle Langevin under Strong Convexity}. LL-03, 2026. Draft report, content revision ce6396dff331. \url{https://odakyildiz.github.io/LangevinLean/docs/pdf/IPLA/FullProof.pdf}.
\end{thebibliography}

\end{document}
